隨著量子計算技術的快速發展，許多大型科技公司近期相繼推出量子計算雲端平台，同時，較小型的科技新創公司與研究團隊也發布了針對含噪聲中等規模量子（NISQ）設備的量子模擬函式庫。值得注意的是，目前大多數量子機器皆採用基於邏輯閘的架構，量子電路即為此類平台的核心基礎。從密碼分析攻擊元件到量子金鑰分發協定，所有量子密碼學演算法最終都必須以電路的形式執行。由於量子計算存在去相干（decoherence）問題，會使儲存於量子硬體上的資訊完整性隨之劣化，因此使用愈多量子邏輯閘將導致電路深度增加，會相應降低運算準確度。因故，最小化電路深度仍是提升此類架構下密碼系統效能的關鍵指標。

在我們的研究中，我們基於 Shi 與 Feng 所提出的貪婪演算法，向對稱式區塊加密演算法的量子電路進行最佳化。我們的方法採雙向搜尋策略，探索所有可能的淺深度 CNOT 電路，以合成量子密碼演算法中所使用的線性可逆元件。我們針對 AES MixColumns 區塊，成功合成出深度為 9 的 CNOT 電路，優於先前的貪婪演算法結果，並證實本方法可推廣應用於其他區塊加密演算法。此外，我們的演算法亦能有效應用於多種 MDS 矩陣，為每一種矩陣合成出低深度電路。再者，我們改良了先前深度導向貪婪演算法研究中的若干關鍵環節，進一步提升其效能表現。最後，我們展示本方法在近期廣受關注的輕量級區塊加密演算法 KLEIN 家族上的實務適用性：透過本方法，我們合成出一個資源效率更高的錯位（out-of-place）MixNibbles 元件，相較於原始設計，CNOT 閘數量減少 64.4\%，電路深度降低 94.2\%，同時仍滿足所需的擴散性質。

\medskip
\textbf{關鍵字：}量⼦電路、量⼦深度、貪婪演算法、區塊加密演算法、AES MixColumns、MDS 矩陣、KLEIN